%==============================================================================
%  1. Introduction
%  写法约定：正文英文、注释中文；中文不得进入正文或 \todo{}（pdflatex 不支持）
%==============================================================================

\section{Introduction}
\label{sec:intro}

% 段落规划（可删）：
%   P1 应用背景与工程动机
%   P2 现有方法及其局限（DMDE 的 CR / F / GMR 被手工方程耦合）
%   P3 LLM 用于优化的相关进展，指出 gap
%   P4 本文贡献（bullet 列表）
%   P5 论文结构

Multi-UAV cooperative path planning requires a set of trajectories that
simultaneously avoids terrain obstacles and radar threat fields while satisfying
kinematic constraints and mission-level task assignment.
\todo{state the essential difficulty: coupled constraints, non-convexity,
expensive evaluation}.
Differential evolution (DE) \citep{storn1997de} has become a popular solver for
such continuous optimization problems owing to its simple structure and strong
robustness.

\todo{Paragraph 2: present the DMDE formulation and explain why the three
search mechanisms -- the crossover rate \CR{}, the scale factor \Ffactor{} and
the extinction rate \GMR{} -- are coupled by hand-designed equations
(formulas~3-9, 3-11 and 3-12), so that they can only be moved together and
cannot respond independently to the observed search state.}

\todo{Paragraph 3: review representative work on using LLMs for optimization
(hyper-parameter control, operator design, heuristic generation) and state the
gap: most approaches treat the LLM as a one-shot generator and lack a closed
loop that feeds the outcome of the previous decision back to the model.}

The main contributions of this work are summarised as follows.

\begin{itemize}
  \item We propose an \llmdmde{} framework that decouples the three search
        mechanisms of DMDE -- \CR{}, \Ffactor{} and \GMR{} -- and delegates
        their coordination to an LLM, while leaving the evolutionary operators
        untouched so that it can be plugged into existing DE implementations.
  \item We design \searchctrl{}, in which the LLM is queried at fixed intervals
        and receives explicit feedback on its previous decision together with
        its fitness, diversity and feasibility consequences, and against a
        fixed-\CR{} shadow population that makes the feedback attributable,
        forming a closed decision loop.
  \item We isolate the effect of the action space itself: a decoupled, a
        coupled (\CR{} only) and a CR-locked (\Ffactor{} and \GMR{} only)
        condition share the same observations, triggers and guards, so that
        only what the model may control differs.
  \item We report a negative result with a mechanistic explanation: on this
        static assignment task an offline constant \CR{} = 0.3 outperforms all
        online variants, and we attribute this to the granularity of the action
        space, the signal-to-noise ratio of a single decision, and the absence
        of a numeric-calibration prior in the pretrained model.
\end{itemize}

The remainder of this paper is organised as follows. \secref{sec:problem}
presents the problem formulation and the \dmde{} baseline; \secref{sec:method}
describes the proposed framework; \secref{sec:experiments} reports and
discusses the experimental results; \secref{sec:conclusion} concludes the paper.
